\section{Privacy}
\label{sec:privacy}

\subsection{Threat Model}

The AI Chain assumes three independent adversaries:

\begin{itemize}[leftmargin=1.2em]
  \item A \emph{network-level} adversary that observes traffic between
        callers, operators, and the chain.
  \item A \emph{provider-level} adversary that operates GPUs and runs
        inference on caller data.
  \item A \emph{chain-level} adversary that observes everything written
        on chain.
\end{itemize}

Plaintext on chain is public to all three. The privacy primitives below
move sensitive material outside the chain entirely or hide it under
encryption that no level of the adversary can decrypt without keys.

\subsection{Public Inference}

The default mode is public inference. The caller sends plaintext input;
the operator returns plaintext output; the receipt commits to both. This
mode is appropriate for non-sensitive workloads (general chat, public
data analysis, code completion on open-source code). The chain records
input\_root and output\_root commitments — Merkle roots over the
plaintext — but the underlying bytes live in CAS and are not stored on
chain.

\subsection{Confidential Inference (F-Chain Delegation)}

For sensitive workloads — medical, financial, defense, and any case
where the input itself is the sensitive asset — the AI Chain delegates
to F-Chain (\S\ref{sec:xchain}) under an activated confidentiality profile.
That profile may use an approved homomorphic-encryption scheme such as CKKS,
confidential hardware, threshold decryption, or a governed composition. Its
exact parameters, supported operations, leakage model, and trust assumptions
must be fixed before the job is admitted. A TEE is optional rather than implied
by homomorphic execution, and no profile may claim that an operator cannot see
plaintext unless the activated construction actually provides that property.

The F-Chain inference path inherits the LP-067 confidential ERC-20
pattern in spirit: the on-chain object commits to a value and an
authority but does not reveal the value. Here the object is a model
query and an output, not a token amount. When a TEE profile is selected, its
quote binds the input ciphertext root, model commitment, output ciphertext root,
job, runtime measurement, and freshness to the declared hardware identity. A
non-TEE profile must supply the proof and key-management bindings specified by
its own policy.

\subsection{Threshold Custody for Model Weights}

Model owners may custody decryption keys for their weights under M-Chain
threshold MPC (\S\ref{sec:xchain}). This addresses two scenarios:

\begin{itemize}[leftmargin=1.2em]
  \item \emph{Selective release}: the model's weights are encrypted at
        rest; an operator that wants to load the model must obtain a
        threshold-signed decryption shard. Custody quorums can be drawn
        to satisfy contractual or jurisdictional requirements.
  \item \emph{Revocable inference}: a model that has been served for a
        time can be retired by retiring the custody key — operators
        cannot reload the weights once the custody quorum withholds the
        shard. Useful for time-bounded research access or for
        deprecating compromised checkpoints.
\end{itemize}

\subsection{Receipt Privacy}

Even with FHE-protected inputs and outputs, the existence of a receipt
is observable on chain. The chain records (caller, operator,
model\_root, block, payment). For workloads where even this metadata
is sensitive, callers may use a stealth-address pattern: the caller is
a one-shot address, the payment is denominated in a confidential AI
balance on F-Chain, and the receipt is encrypted to the caller's
session key. The on-chain footprint becomes a generic ``some address
paid some operator for some model.''

\subsection{Privacy Layering Summary}

\begin{table}[h]
\centering
\small
\begin{tabular}{lccc}
\toprule
\textbf{Mode} & \textbf{Input} & \textbf{Output} & \textbf{Identity} \\
\midrule
Public               & plain & plain & visible \\
Confidential (F)     & FHE   & FHE   & visible \\
Custody (M)          & plain & plain & visible \\
Custody + Confidential & FHE & FHE & visible \\
Confidential + Stealth & FHE & FHE & hidden \\
\bottomrule
\end{tabular}
\caption{Privacy modes available on AI Chain. Each row composes
independent primitives; the chain validates the receipt regardless of
which combination the caller chose.}
\end{table}

\subsection{What Is Not Hidden}

The chain does not hide block height, the validator set, the existence
of a receipt, the finality proof, or the operator's stake state.
These are network invariants and must be public for the consensus
itself to function. Anything tied to the workload (input, output,
model weights when custodied, optionally the caller identity) can be
hidden behind the primitives above.
